% Author: Keshav
\begin{abstract}
We link noisy business records from two secondary sources to a reference registry of US and Indian
businesses whose names and addresses are written in Latin, Devanagari, Malayalam and other Indic scripts.
Every model is evaluated with the competition metric, macro \fhalf{} per registry entity, on a held-out
split of four whole states (21,529 test entities). A shared normalisation step romanises non-Latin text
and applies a transliteration lexicon learned from the training links. Candidates come from character
3-gram TF-IDF and a fine-tuned multilingual-e5 bi-encoder, fused by reciprocal rank fusion (top 5,
ceiling 0.99957). On the same candidate pairs we compare matchers from three model families: a LightGBM
model over edit-distance and retrieval features (0.9908), Siamese character RNNs (best 0.9834), and
transformer cross-encoders: zero-shot, fully fine-tuned, LoRA, a Mixture-of-Experts head,
IndicBERTv2, mDeBERTa-v3 (best single model, 0.9940) and Ditto (0.9876 at 160k pairs). A gated Mixture of Experts over nine
matchers reaches 0.9948. Four findings stand out. Task fine-tuning matters more than the backbone
(zero-shot 0.9167 vs.\ fine-tuned 0.9912). The classical model is competitive because 83\% of its gain
comes from the fine-tuned retriever. An Indic-specific backbone loses to a multilingual reranker even on
Indic scripts. The fusion gain comes from combining two complementary experts, not from per-pair routing.
\end{abstract}
